\subsection{OpenCL -- Open Computing Language}
\label{sec:frameworks:opencl}

\begin{quotebox}{\href{https://www.khronos.org/opencl}{Khronos OpenCL}}
    OpenCL™ (Open Computing Language) is an open, royalty-free standard for cross-platform, parallel programming of diverse accelerators found in supercomputers, cloud servers, personal computers, mobile devices and embedded platforms.
\end{quotebox}

\acrfull{opencl} was the first standardized programming framework with which programmers were able to target \glspl{cpu} and \glspl{gpu} (and many more hardware platforms) using a single source code. 
The first \gls{opencl} Khronos working group was created in 2008~\cite{khronos_opencl_working_group_khronos_2008}, shortly after the initial \gls{cuda} release, and the first official \gls{opencl} specification was released one year later in 2009~\cite{munshi_opencl_2009}. 
The current \gls{opencl} version 3.0 revision 17 was published in October 2024~\cite{khronos_opencl_working_group_opencl_2024}. 
As was the case with \gls{openmp}, since \gls{opencl} is a standard, there are many different implementations of \gls{opencl} targeting different hardware platforms. 
These implementations include, among others, implementations of the \gls{opencl} 3.0 standard from Intel, Google, ARM, NVIDIA, or PoCL. 
Many more implementations exist for the \gls{opencl} 2.0 standard\footnote{\url{https://www.khronos.org/conformance/adopters/conformant-products/opencl} (visited: 2025-03-21)}.

\begin{listing}
    \begin{moderncode}[fontsize=\small]{cpp}
const char* kernel_source = R"KS(
__kernel void saxpy(const double alpha, __global double *X, __global double *X) {
    const int idx = get_global_id(0);
    Y[idx] = alpha * X[idx] + Y[idx];
}
)KS";

cl_platform_id platform;
clGetPlatformIDs(1, &platform, nullptr);
cl_device_id device;
clGetDeviceIDs(platform, CL_DEVICE_TYPE_GPU, 1, &device, nullptr);
cl_context context = clCreateContext(nullptr, 1, &device, nullptr, nullptr, nullptr);
cl_command_queue queue = clCreateCommandQueue(context, device, 0, nullptr);

cl_mem d_X = clCreateBuffer(context, CL_MEM_READ_ONLY | CL_MEM_COPY_HOST_PTR, N * sizeof(double), X, nullptr);
cl_mem d_Y = clCreateBuffer(context, CL_MEM_READ_ONLY | CL_MEM_COPY_HOST_PTR, N * sizeof(double), Y, nullptr);

cl_program program = clCreateProgramWithSource(context, 1, &kernel_source, nullptr, nullptr);
clBuildProgram(program, 1, &device, nullptr, nullptr, nullptr);
cl_kernel kernel = clCreateKernel(program, "saxpy", nullptr);

clSetKernelArg(kernel, 0, sizeof(double), &alpha);
clSetKernelArg(kernel, 1, sizeof(cl_mem), &d_X);
clSetKernelArg(kernel, 2, sizeof(cl_mem), &d_Y);

size_t global_size = N;
clEnqueueNDRangeKernel(queue, kernel, 1, nullptr, &global_size, nullptr, 0, nullptr, nullptr);
clEnqueueReadBuffer(queue, d_C, CL_TRUE, 0, N * sizeof(double), C, 0, nullptr, nullptr);

clReleaseMemObject(d_X); 
clReleaseMemObject(d_Y);
clReleaseKernel(kernel); 
clReleaseProgram(program);
clReleaseCommandQueue(queue); 
clReleaseContext(context);
    \end{moderncode}
    \caption{%
    A simple DAXPY example using \gls{opencl} including all necessary compilation, memory, and kernel launch operations. }
    \label{code:opencl_example}
\end{listing}

\gls{opencl} follows the same host-to-device offloading model as \gls{cuda} including the mentioned general workflow. 
However, as can be seen in \autoref{code:opencl_example}, \gls{opencl} is noticeably more verbose and involved than \gls{cuda}. 
\gls{opencl} allows the programmer to dynamically specify to which device the code should be offloaded at runtime. 
For that, the \code{cl_command_queue} (lines 8-13) has to be created, which is later used to enqueue the actual compute kernels. 
Additionally, in \gls{opencl}, it is necessary to compile the kernel source, given as a simple C-style string, by hand, as well as setting the actual kernel arguments (lines 18-24). 
This approach gives \gls{opencl} great flexibility since the compute kernels can easily be assembled and modified at runtime. 
At the same time, however, the verbosity of \gls{opencl} also opens the doors for many more pitfalls and potential bugs. 
Also, while \gls{opencl} is not as high-level as \gls{openmp}, not all vendor-specific optimizations are possible since one of the main goals of \gls{opencl} is to support as many platforms as possible. 

\todo{different parallelism strategies?}

%%% ADVANTAGES AND DISADVANTAGES
\AdvantagesAndDisadvantages{opencl}{
    \item standardized
    \item supports a wide variety of hardware platforms
}{
    \item kernels must be compiled by hand
    \item error-prone
    \item some low-level optimizations not possible
}
